\newpage
\section{Implementacja}

\subsection{Trening modelu}

Najważniejszym elementem implementacji jest skrypt \texttt{train.py}, który odpowiada za:

\begin{itemize}
    \item wczytanie konfiguracji z pliku \texttt{config.yaml},
    \item przygotowanie modelu,
    \item trening i walidację,
    \item zapis najlepszych wag,
    \item wznowienie treningu z checkpointu,
    \item zapis historii uczenia i podsumowania do pliku JSON.
\end{itemize}

W trakcie pracy dodano także mechanizmy praktyczne, które okazały się istotne dla dłuższych eksperymentów:

\begin{itemize}
    \item \texttt{early stopping},
    \item checkpointy \texttt{best\_model.pt}, \texttt{last\_checkpoint.pt} i \texttt{interrupted\_checkpoint.pt},
    \item obsługę treningu na GPU z wykorzystaniem \texttt{AMP},
    \item automatyczny dobór rozmiaru batcha,
    \item lepsze ustawienia loadera pod Windows i CUDA.
\end{itemize}

\subsection{Metryki}

Podczas ewaluacji modelu raportowane są:

\begin{itemize}
    \item \texttt{Accuracy},
    \item \texttt{Precision Macro},
    \item \texttt{Recall Macro},
    \item \texttt{F1 Macro},
    \item \texttt{ROC-AUC}.
\end{itemize}

Metryki te pozwalają ocenić model nie tylko pod względem skuteczności ogólnej, ale też pod kątem jakości rozróżniania klas.

\subsection{Wyniki etapu bazowego}

Wcześniejszy model bazowy osiągnął na zbiorze testowym:

\begin{itemize}
    \item Accuracy: 0.9339,
    \item F1 Macro: 0.9339,
    \item ROC-AUC: 0.9831.
\end{itemize}

Następnie przeprowadzono drugi trening z mocniejszymi augmentacjami odpornościowymi. Wyniki tego wariantu były następujące:

\begin{itemize}
    \item Accuracy: 0.9268,
    \item Precision Macro: 0.9273,
    \item Recall Macro: 0.9267,
    \item F1 Macro: 0.9267,
    \item ROC-AUC: 0.9804.
\end{itemize}

Wariant ten okazał się minimalnie słabszy na czystym zbiorze testowym, ale lepiej przygotowany do dalszych analiz odporności.

\subsection{Etap twarzowy}

W drugim etapie projektu przygotowano osobną ścieżkę analizy portretów i cropów twarzy. W praktyce oznacza to:

\begin{itemize}
    \item automatyczne wykrycie twarzy klasyfikatorem Haar Cascade,
    \item zapis cropu w stylu \texttt{face} albo \texttt{portrait},
    \item trening osobnego modelu twarzowego niezależnego od klasyfikatora globalnego,
    \item możliwość porównania modelu globalnego i modelu twarzowego na tym samym zbiorze portretów.
\end{itemize}

Wariant twarzowy został oparty o architekturę \texttt{ConvNeXt Tiny}. Dla curated miksu wieloźródłowego uzyskano bardzo dobre wyniki na danych zbliżonych do treningowych:

\begin{itemize}
    \item Accuracy: 0.9763,
    \item Precision Macro: 0.9765,
    \item Recall Macro: 0.9763,
    \item F1 Macro: 0.9763,
    \item ROC-AUC: 0.9982.
\end{itemize}

Pokazało to, że dla klasycznych danych portretowych model twarzowy może być bardzo skuteczny, ale nie przesądza jeszcze o jego odporności na nowe generatory obrazów.

\subsection{Test odporności}

Dla aktualnego modelu przeprowadzono osobny test odporności na spadek jakości obrazu. Sprawdzono między innymi przypadki:

\begin{itemize}
    \item lekkiej i silnej kompresji JPEG,
    \item lekkiego i silnego rozmycia,
    \item lekkiego i silnego downscale,
    \item profilu mieszanego \texttt{mixed\_quality}.
\end{itemize}

Najważniejsze obserwacje były następujące:

\begin{itemize}
    \item lekkie zaburzenia obniżały skuteczność tylko nieznacznie,
    \item najsilniejszy spadek pojawiał się dla \texttt{jpeg\_extreme},
    \item przypadki mieszane obniżały skuteczność wyraźnie bardziej niż pojedyncza lekka degradacja.
\end{itemize}

Pokazuje to, że model jest już w pewnym stopniu odporny na pogorszenie jakości, ale nadal pozostaje wrażliwy na obrazy bardzo mocno skompresowane.

\subsection{Benchmark OOD dla nowych generatorów}

Aby sprawdzić, czy modele generalizują poza główny rozkład treningowy, przygotowano mały benchmark OOD złożony z:

\begin{itemize}
    \item 20 obrazów wygenerowanych przez GPT,
    \item 20 obrazów wygenerowanych przez Gemini / Nano2,
    \item 40 prawdziwych zdjęć portretowych dobranych tak, aby stylistycznie przypominały obrazy syntetyczne.
\end{itemize}

Benchmark ten nie był częścią głównego zbioru treningowego. Jego celem było możliwie uczciwe sprawdzenie, jak modele zachowują się na nowszych generatorach i bardziej współczesnych portretach. W praktyce był to więc mały test zewnętrzny, pozwalający sprawdzić nie tylko skuteczność klasyfikacji, ale także odporność na przesunięcie rozkładu danych.

Na tym benchmarku model globalny uzyskał:

\begin{itemize}
    \item Accuracy: 0.7125,
    \item Precision Macro: 0.7934,
    \item Recall Macro: 0.7125,
    \item F1 Macro: 0.6912,
    \item ROC-AUC: 0.8413.
\end{itemize}

Macierz pomyłek pokazała istotną asymetrię:

\begin{itemize}
    \item 39 z 40 obrazów \texttt{fake} zostało wykrytych poprawnie,
    \item tylko 18 z 40 obrazów \texttt{real} zostało rozpoznanych poprawnie,
    \item aż 22 prawdziwe zdjęcia zostały oznaczone jako \texttt{fake}.
\end{itemize}

Wniosek z tego etapu jest ważny: model globalny nie zawodził głównie na nowych obrazach syntetycznych, lecz na nowych realistycznych portretach, dla których generował zbyt wiele fałszywych alarmów.

\subsection{Tabela zbiorcza wyników}

Najważniejsze wyniki eksperymentów zestawiono w tabeli \ref{tab:wyniki-zbiorcze}. Pozwala ona porównać modele nie tylko na danych zbliżonych do treningowych, ale również na małym benchmarku OOD oraz po krótkiej adaptacji modelu twarzowego.

\begin{table}[H]
\centering
\caption{Zbiorcze wyniki najważniejszych eksperymentów}
\label{tab:wyniki-zbiorcze}
\begin{tabular}{|p{4.3cm}|c|c|c|}
\hline
\textbf{Wariant} & \textbf{Accuracy} & \textbf{F1 Macro} & \textbf{ROC-AUC} \\
\hline
Model globalny -- baseline test & 0.9339 & 0.9339 & 0.9831 \\
\hline
Model globalny -- wariant robust & 0.9268 & 0.9267 & 0.9804 \\
\hline
Model twarzowy \texttt{ConvNeXt Tiny} -- curated test & 0.9763 & 0.9763 & 0.9982 \\
\hline
Model globalny -- benchmark OOD & 0.7125 & 0.6912 & 0.8413 \\
\hline
Model twarzowy przed adaptacją -- OOD (wykryte twarze) & ok. 0.51 & ok. 0.49 & ok. 0.52 \\
\hline
Model twarzowy po adaptacji -- test wewnętrzny & 0.8548 & 0.8548 & 0.9469 \\
\hline
\end{tabular}
\end{table}

\subsection{Analiza jakościowa błędów}

Same metryki nie pokazują jeszcze, \emph{jakiego rodzaju} obrazy są problematyczne. Dlatego wykonano również analizę jakościową przypadków błędnych klasyfikacji.

Na rysunku \ref{fig:ood-false-positive} pokazano przykład fałszywego alarmu modelu globalnego na prawdziwym portrecie. Tego typu przypadki były częste w benchmarku OOD i sugerowały, że model nauczył się częściowo traktować estetyczne, czyste portrety jako obrazy podejrzane.

\begin{figure}[H]
\centering
\includegraphics[width=0.42\textwidth]{../reports/ood_pilot_global/test/examples/false_positive/04__gt_real__pred_fake__real_matched_real__0002__pexels-a-darmel-8158953_jpg.jpg}
\caption{Przykład false positive na benchmarku OOD: prawdziwy portret sklasyfikowany jako \texttt{fake}}
\label{fig:ood-false-positive}
\end{figure}

Na rysunku \ref{fig:ood-false-negative} pokazano odwrotny przypadek, czyli obraz syntetyczny wygenerowany przez GPT, który został oznaczony jako \texttt{real}. Takie próbki były rzadsze niż false positive modelu globalnego, ale bardzo istotne z punktu widzenia wiarygodności systemu.

\begin{figure}[H]
\centering
\includegraphics[width=0.42\textwidth]{../reports/ood_pilot_global/test/examples/false_negative/01__gt_fake__pred_real__fake_gpt__0016__gpt__p017__01_png.png}
\caption{Przykład false negative na benchmarku OOD: obraz syntetyczny sklasyfikowany jako \texttt{real}}
\label{fig:ood-false-negative}
\end{figure}

Takie przykłady potwierdzają, że problem generalizacji nie dotyczy wyłącznie pojedynczych metryk, ale ma również wymiar jakościowy: modele potrafią mylić się na wizualnie przekonujących portretach nawet wtedy, gdy na danych wewnętrznych osiągają bardzo wysokie wyniki.

\subsection{Wyniki modelu twarzowego na nowych portretach}

Porównanie wcześniejszych checkpointów twarzowych z benchmarkiem OOD pokazało, że sam dobry wynik na danych wewnętrznych nie wystarcza do poprawnej generalizacji na nowsze obrazy. Przed adaptacją:

\begin{itemize}
    \item wariant \texttt{faces\_hard\_adapt} uzyskał Accuracy około 0.48 na podzbiorze z wykrytą twarzą,
    \item wariant \texttt{faces\_convnext\_v2} uzyskał Accuracy około 0.51 na tym samym benchmarku.
\end{itemize}

Oznacza to, że pierwotna wersja modelu twarzowego nie generalizowała poprawnie na najnowsze fotorealistyczne obrazy GPT / Gemini.

W dalszym kroku wykonano mały fine-tuning modelu \texttt{faces\_convnext\_v2} na cropach twarzy przygotowanych z nowych obrazów. Na wewnętrznym teście adaptacyjnym uzyskano:

\begin{itemize}
    \item Accuracy: 0.8548,
    \item Precision Macro: 0.8552,
    \item Recall Macro: 0.8548,
    \item F1 Macro: 0.8548,
    \item ROC-AUC: 0.9469.
\end{itemize}

Wynik ten należy jednak interpretować ostrożnie, ponieważ dotyczy niewielkiego zbioru adaptacyjnego, a ręczne testy w interfejsie nadal pokazywały trudne przypadki, w których nowe obrazy syntetyczne były klasyfikowane jako \texttt{real}. Oznacza to, że adaptacja poprawiła zachowanie modelu, ale nie rozwiązała jeszcze problemu w sposób pełny.

\subsection{Interfejs demonstracyjny}

W celu zwiększenia użyteczności projektu przygotowano lokalne demo webowe w \texttt{Gradio}. Aplikacja umożliwia:

\begin{itemize}
    \item załadowanie pojedynczego obrazu,
    \item wyświetlenie predykcji \texttt{real/fake},
    \item pokazanie pewności modelu,
    \item wygenerowanie mapy Grad-CAM,
    \item pokazanie ostrzeżenia przy niskiej jakości obrazu,
    \item porównanie predykcji modelu globalnego i modelu twarzowego dla portretów.
\end{itemize}

Interfejs ten ma charakter demonstracyjny i stanowi przydatne narzędzie do ręcznej analizy zachowania modelu.
